\documentclass[11pt,a4paper]{article}
\usepackage[margin=1in]{geometry}
\usepackage{amsmath,amssymb,amsthm,mathtools}
\usepackage{enumitem}
\usepackage{hyperref}
\usepackage{microtype}
\hypersetup{colorlinks=true,linkcolor=blue}
\newcommand{\Exp}{\operatorname{Exp}}
\newcommand{\Log}{\operatorname{Log}}
\newcommand{\SO}{\mathrm{SO}}
\newcommand{\R}{\mathbb{R}}
\newcommand{\defeq}{\mathrel{\mathop:}=}
\title{Chapter 4: IMU Pre-integration\\\large Derivation of the Three Pre-integrated Increments}
\author{}
\date{}
\usepackage{graphicx}
\begin{document}
\maketitle
\tableofcontents

\section{Notation and assumptions}
We derive the three pre-integrated increments corresponding to the textbook's equations (4.7), starting from the discrete-time IMU kinematics (4.4) and their accumulated form (4.6). Unlike the textbook's constant sampling interval $\Delta t$, we allow unequal sample intervals:
\begin{equation}
\boxed{\Delta t_k=t_{k+1}-t_k,\qquad \Delta t_{ij}=t_j-t_i=\sum_{k=i}^{j-1}\Delta t_k.}
\end{equation}
We assume $R_k\in\SO(3)$ maps body-frame vectors to world-frame vectors, $p_k,v_k,g\in\R^3$ are world-frame quantities, gravity $g$ is constant over the interval, and the bias/noise terms are represented by $b_{g,k},b_{a,k},\eta_{gd,k},\eta_{ad,k}$. The IMU measurements are $\tilde\omega_k$ and $\tilde a_k$. We use a zero-order-hold approximation over each sample interval for the corrected acceleration and angular velocity. All products of rotations are \emph{time ordered}.


\subsection{Continuous-time kinematics (textbook Section 4.1.1)}
The starting point is the continuous-time rigid-body kinematics,
where $a(t)$ is the \emph{world-frame total acceleration} and
$\omega(t)$ is the body-frame angular velocity:
\begin{subequations}\label{eq:continuous_kinematics}
\begin{align}
\dot R(t)&=R(t)\omega(t)^\wedge,\\
\dot v(t)&=a(t),\\
\dot p(t)&=v(t).
\end{align}
\end{subequations}
On the interval $[t_k,t_{k+1}]$, let $\Delta t_k=t_{k+1}-t_k$.
If $\omega_k$ is held constant, solving the rotation differential
 equation gives
\begin{equation}
R_{k+1}=R_k\Exp(\omega_k\Delta t_k).
\end{equation}
With constant world-frame acceleration $a_k$ over the step,
 integration of $\dot v=a_k$ gives
\begin{equation}
v_{k+1}=v_k+a_k\Delta t_k.
\end{equation}
Since $v(t_k+\tau)=v_k+a_k\tau$, integrating $\dot p=v$ gives
\begin{align}
p_{k+1}&=p_k+\int_0^{\Delta t_k}(v_k+a_k\tau)\,d\tau\notag\\
&=p_k+v_k\Delta t_k+\tfrac12a_k\Delta t_k^2.
\end{align}
These are the textbook's single-step discrete kinematics.
For the position and velocity steps, holding $a_k$ constant is an
approximation: during a real IMU interval the body may rotate, so
$R(t)$ and hence the world-frame acceleration can vary.

\subsection{Continuous-time IMU measurement model and gravity}
Before discretizing the sensor noise, use the continuous-time IMU
measurement model of textbook Section 4.1.1. The gyroscope measures
body-frame angular velocity, and the accelerometer measures body-frame
\emph{specific force} (the acceleration excluding gravity):
\begin{subequations}\label{eq:imu_model}
\begin{align}
\tilde\omega(t)&=\omega(t)+b_g(t)+\eta_g(t),\\
\tilde a(t)&=R(t)^T\bigl(a(t)-g\bigr)+b_a(t)+\eta_a(t).
\end{align}
\end{subequations}
Here $\eta_g(t)$ and $\eta_a(t)$ denote continuous-time white
measurement noise, and $b_g(t),b_a(t)$ are time-dependent biases.
Solving for the quantities appearing in the kinematics gives
\begin{align}
\omega(t)&=\tilde\omega(t)-b_g(t)-\eta_g(t),\\
a(t)&=g+R(t)\bigl(\tilde a(t)-b_a(t)-\eta_a(t)\bigr).
\end{align}
Thus gravity appears in the world-frame acceleration even though the
accelerometer measures specific force in the body frame.

\subsection{Discrete-time noise convention}
To apply the continuous-time measurement model over
$[t_k,t_{k+1}]$, use the sample interval $\Delta t_k$ and the
interval-averaged discrete noises $\eta_{gd,k}$ and $\eta_{ad,k}$.
If $Q_{g,c}$ and $Q_{a,c}$ are the continuous-time white-noise
intensity matrices, their discrete sample covariances are
\begin{subequations}\label{eq:noise_discretization_intro}
\begin{align}
\operatorname{Cov}(\eta_{gd,k})&=\frac{Q_{g,c}}{\Delta t_k},\\
\operatorname{Cov}(\eta_{ad,k})&=\frac{Q_{a,c}}{\Delta t_k}.
\end{align}
\end{subequations}
These are the discrete measurement-noise samples appearing in the
single-step equations below. The factors $1/\Delta t_k$ refer to
sample-average noise covariance; the covariance of the integrated
noise increment instead scales with $\Delta t_k$.

\subsection{Single-step discrete IMU propagation}
Evaluating the continuous-time measurement model at each sample and
using the discrete noise convention above, substitute the corrected
angular velocity and world-frame acceleration into the single-step
kinematics. This gives the three equations used throughout the derivation:
\begin{subequations}\label{eq:step}
\begin{align}
R_{k+1}&=R_k\Exp\!\left((\tilde\omega_k-b_{g,k}-\eta_{gd,k})\Delta t_k\right),\\
v_{k+1}&=v_k+g\Delta t_k
 +R_k(\tilde a_k-b_{a,k}-\eta_{ad,k})\Delta t_k,\\
p_{k+1}&=p_k+v_k\Delta t_k+\tfrac12g\Delta t_k^2\notag\\
&\quad+\tfrac12R_k(\tilde a_k-b_{a,k}-\eta_{ad,k})\Delta t_k^2.
\end{align}
\end{subequations}
The gravity terms follow directly from the world-frame acceleration
$a_k=g+R_k(\tilde a_k-b_{a,k}-\eta_{ad,k})$.
In particular, the factor $\tfrac12$ in the position equation comes
from integrating velocity over one sample interval.

\section{Rotation increment $\Delta R_{ij}$}
Repeated substitution into \eqref{eq:step}a gives
\begin{align*}
R_{i+1}&=R_i\Exp((\tilde\omega_i-b_{g,i}-\eta_{gd,i})\Delta t_i),\\
R_j&=R_i\prod_{k=i}^{j-1}
\Exp((\tilde\omega_k-b_{g,k}-\eta_{gd,k})\Delta t_k).
\end{align*}
Left-multiplying by $R_i^T$ and using $R_i^TR_i=I_3$ yields
\begin{equation}\label{eq:dR}
\boxed{\Delta R_{ij}\defeq R_i^TR_j=\prod_{k=i}^{j-1}\Exp\!\left((\tilde\omega_k-b_{g,k}-\eta_{gd,k})\Delta t_k\right).}
\end{equation}
The product is ordered from the earliest sample $i$ on the left to the latest sample $j-1$ on the right; its factors cannot generally be reordered. In particular, $\Delta R_{ii}=I_3$ and
\begin{equation}\label{eq:Rk}
\Delta R_{ik}=R_i^TR_k,\qquad R_k=R_i\Delta R_{ik}.
\end{equation}
The right-hand side of \eqref{eq:dR} depends on IMU readings, biases, and noise, but not on the starting attitude $R_i$.

\section{Velocity increment $\Delta v_{ij}$}
Summing \eqref{eq:step}b from $k=i$ to $j-1$ gives a telescoping sum on the left:
\begin{align}
v_j-v_i
&=\sum_{k=i}^{j-1}(v_{k+1}-v_k)\notag\\
&=\sum_{k=i}^{j-1}g\Delta t_k+\sum_{k=i}^{j-1}R_k(\tilde a_k-b_{a,k}-\eta_{ad,k})\Delta t_k\notag\\
&=g\Delta t_{ij}+\sum_{k=i}^{j-1}R_k(\tilde a_k-b_{a,k}-\eta_{ad,k})\Delta t_k.
\end{align}
Substitute $R_k=R_i\Delta R_{ik}$ from \eqref{eq:Rk}:
\begin{align}
v_j-v_i-g\Delta t_{ij}
&=\sum_{k=i}^{j-1}R_i\Delta R_{ik}(\tilde a_k-b_{a,k}-\eta_{ad,k})\Delta t_k\notag\\
&=R_i\sum_{k=i}^{j-1}\Delta R_{ik}(\tilde a_k-b_{a,k}-\eta_{ad,k})\Delta t_k.
\end{align}
Left-multiply by $R_i^T$:
\begin{equation}\label{eq:dv}
\boxed{\begin{aligned}
\Delta v_{ij}&\defeq R_i^T(v_j-v_i-g\Delta t_{ij})\\
&=\sum_{k=i}^{j-1}\Delta R_{ik}(\tilde a_k-b_{a,k}-\eta_{ad,k})\Delta t_k.
\end{aligned}}
\end{equation}
In particular, for any intermediate $k$,
\begin{equation}\label{eq:vk}
\Delta v_{ik}=R_i^T(v_k-v_i-g\Delta t_{ik})
\quad\Longrightarrow\quad
v_k=v_i+g\Delta t_{ik}+R_i\Delta v_{ik}.
\end{equation}
Thus $\Delta v_{ij}$ is \emph{not} simply $v_j-v_i$: it removes gravity's contribution and expresses the result in the starting body frame.

\section{Position increment $\Delta p_{ij}$}
\subsection{Accumulate the position equation}
Summing \eqref{eq:step}c gives
\begin{equation}\label{eq:pacc}
p_j-p_i
=\sum_{k=i}^{j-1}v_k\Delta t_k
+\tfrac12g\sum_{k=i}^{j-1}\Delta t_k^2
+\tfrac12\sum_{k=i}^{j-1}R_k(\tilde a_k-b_{a,k}-\eta_{ad,k})\Delta t_k^2.
\end{equation}
Unlike the velocity case, the first sum contains intermediate velocities $v_k$. We eliminate them using \eqref{eq:vk}.

\subsection{Substitute the intermediate velocity}
\begin{align}
\sum_{k=i}^{j-1}v_k\Delta t_k
&=\sum_{k=i}^{j-1}(v_i+g\Delta t_{ik}+R_i\Delta v_{ik})\Delta t_k\notag\\
&=v_i\Delta t_{ij}+g\sum_{k=i}^{j-1}\Delta t_{ik}\Delta t_k
+R_i\sum_{k=i}^{j-1}\Delta v_{ik}\Delta t_k.
\end{align}
Insert this into \eqref{eq:pacc}:
\begin{align}
p_j-p_i
&=v_i\Delta t_{ij}
+g\left[\sum_{k=i}^{j-1}\Delta t_{ik}\Delta t_k+\tfrac12\sum_{k=i}^{j-1}\Delta t_k^2\right]\notag\\
&\quad+R_i\sum_{k=i}^{j-1}\Delta v_{ik}\Delta t_k
+\tfrac12\sum_{k=i}^{j-1}R_k(\tilde a_k-b_{a,k}-\eta_{ad,k})\Delta t_k^2.
\label{eq:pintermediate}
\end{align}

\subsection{Combine gravity terms using a telescoping sum}
Since $\Delta t_{i,k+1}=\Delta t_{ik}+\Delta t_k$,
\begin{align}
\Delta t_{i,k+1}^2-\Delta t_{ik}^2
&=(\Delta t_{ik}+\Delta t_k)^2-\Delta t_{ik}^2\notag\\
&=2\Delta t_{ik}\Delta t_k+\Delta t_k^2.
\end{align}
Therefore
\begin{equation}
\Delta t_{ik}\Delta t_k+\tfrac12\Delta t_k^2
=\tfrac12(\Delta t_{i,k+1}^2-\Delta t_{ik}^2).
\end{equation}
Sum over all intervals:
\begin{align}
\sum_{k=i}^{j-1}\left(\Delta t_{ik}\Delta t_k+\tfrac12\Delta t_k^2\right)
&=\tfrac12\sum_{k=i}^{j-1}(\Delta t_{i,k+1}^2-\Delta t_{ik}^2)\notag\\
&=\tfrac12(\Delta t_{ij}^2-\Delta t_{ii}^2)\notag\\
&=\boxed{\tfrac12\Delta t_{ij}^2},\qquad \Delta t_{ii}=0.
\end{align}
This identity holds even when the $\Delta t_k$ are unequal.

\subsection{Factor out the initial rotation}
Using $R_k=R_i\Delta R_{ik}$ in \eqref{eq:pintermediate}, and combining the gravity terms, yields
\begin{align}
p_j-p_i
&=v_i\Delta t_{ij}+\tfrac12g\Delta t_{ij}^2\notag\\
&\quad+R_i\sum_{k=i}^{j-1}\left[\Delta v_{ik}\Delta t_k+\tfrac12\Delta R_{ik}(\tilde a_k-b_{a,k}-\eta_{ad,k})\Delta t_k^2\right].
\end{align}
Move the initial-state terms to the left and multiply by $R_i^T$:
\begin{equation}\label{eq:dp}
\boxed{\begin{aligned}
\Delta p_{ij}&\defeq R_i^T\left(p_j-p_i-v_i\Delta t_{ij}-\tfrac12g\Delta t_{ij}^2\right)\\
&=\sum_{k=i}^{j-1}\left[\Delta v_{ik}\Delta t_k+\tfrac12\Delta R_{ik}(\tilde a_k-b_{a,k}-\eta_{ad,k})\Delta t_k^2\right].
\end{aligned}}
\end{equation}
Here $\Delta v_{ik}$ is the \emph{pre-integrated} velocity increment from $i$ to $k$, not the world-frame velocity $v_k$. Its use removes dependence on the initial velocity $v_i$ and initial rotation $R_i$ from the sum.

\section{Summary of the three increments}
Recall $\Delta t_{ij}=\sum_{k=i}^{j-1}\Delta t_k$. The pre-integrated quantities are
\begin{subequations}\label{eq:summary}
\begin{align}
\boxed{\Delta R_{ij}}&=R_i^TR_j
=\prod_{k=i}^{j-1}\Exp\!\left((\tilde\omega_k-b_{g,k}-\eta_{gd,k})\Delta t_k\right),\\
\boxed{\Delta v_{ij}}&=R_i^T(v_j-v_i-g\Delta t_{ij})
=\sum_{k=i}^{j-1}\Delta R_{ik}(\tilde a_k-b_{a,k}-\eta_{ad,k})\Delta t_k,\\
\boxed{\Delta p_{ij}}&=R_i^T\left(p_j-p_i-v_i\Delta t_{ij}-\tfrac12g\Delta t_{ij}^2\right)
=\sum_{k=i}^{j-1}\left(\Delta v_{ik}\Delta t_k+\tfrac12\Delta R_{ik}(\tilde a_k-b_{a,k}-\eta_{ad,k})\Delta t_k^2\right).
\end{align}
\end{subequations}
All three right-hand expressions depend on IMU measurements, biases, noise, and sample intervals, not on the absolute initial state $(R_i,v_i,p_i)$. The bias dependence remains; later pre-integration bias corrections address changes in the bias estimates.

\section{Interpretation and implementation notes}
\begin{enumerate}[leftmargin=*]
\item $\Delta R_{ij}$ is a relative rotation. $\Delta v_{ij}$ and $\Delta p_{ij}$ are gravity-compensated increments expressed in the frame at $i$; they are not simply $v_j-v_i$ and $p_j-p_i$.
\item Initialization at $k=i$ is $\Delta R_{ii}=I_3$, $\Delta v_{ii}=0$, and $\Delta p_{ii}=0$.
\item For unequal sample intervals, replace every single-step $\Delta t$ by $\Delta t_k$, while the accumulated duration is $\Delta t_{ij}$.
\item These are the \emph{noise-containing} relative increments of the textbook's definition. Later, the chapter separates noisy quantities from pre-integrated measurements denoted with a tilde and derives their covariance.
\end{enumerate}

\section{Pre-integration measurement model through the textbook's (4.18)}
\subsection{Fixed-bias assumption and noise convention}
For the pre-integration interval $[t_i,t_j]$, we assume that the
gyroscope and accelerometer biases remain constant and are approximated
by their values at the initial time $t_i$:
\begin{equation}
 b_{g,k}\approx b_{g,i},\qquad b_{a,k}\approx b_{a,i}.
\end{equation}
Changes in the bias estimates are handled using the bias Jacobians
derived in the subsequent Bias Update section. We retain only the
following frequently used right-Jacobian notation:
\begin{equation}
J_{r,k}\defeq J_r\!\left((\tilde\omega_k-b_{g,i})\Delta t_k\right).
\end{equation}
Here $J_r$ denotes the right Jacobian of $\SO(3)$.
The discrete noise samples are $\eta_{gd,k}$ and $\eta_{ad,k}$.
The equations below retain only first-order noise terms and neglect
products of two noise perturbations.

\subsection{Separating rotation measurement and noise}
The right-Jacobian first-order identity is
\begin{equation}
\Exp(\theta+\epsilon)
\simeq\Exp(\theta)\Exp(J_r(\theta)\epsilon).
\end{equation}
Setting $\epsilon=-\eta_{gd,k}\Delta t_k$ gives
\begin{equation}
\Exp((\tilde\omega_k-b_{g,i}-\eta_{gd,k})\Delta t_k)
\simeq
\Exp((\tilde\omega_k-b_{g,i})\Delta t_k)
\Exp(-J_{r,k}\eta_{gd,k}\Delta t_k).
\end{equation}
Define the noise-free computed measurement
\begin{equation}\label{eq:tildeR}
\boxed{\Delta\tilde R_{ij}\defeq
\prod_{k=i}^{j-1}\Exp((\tilde\omega_k-b_{g,i})\Delta t_k),\qquad
\Delta\tilde R_{ii}=I_3.}
\end{equation}
For any $B\in\SO(3)$, the adjoint identity is
\begin{equation}\label{eq:adjmove}
\Exp(\xi)B=B\Exp(B^T\xi).
\end{equation}
For example, the first two increments become
\begin{align}
\Delta R_{i,i+2}
&\simeq \Delta\tilde R_{i,i+1}
 \Exp(-J_{r,i}\eta_{gd,i}\Delta t_i)
 \Delta\tilde R_{i+1,i+2}
 \Exp(-J_{r,i+1}\eta_{gd,i+1}\Delta t_{i+1})\notag\\
&=\Delta\tilde R_{i,i+2}
 \Exp(-\Delta\tilde R_{i+1,i+2}^{T}J_{r,i}\eta_{gd,i}\Delta t_i)
 \Exp(-J_{r,i+1}\eta_{gd,i+1}\Delta t_{i+1}).
\end{align}
In particular, the noise accompanying the interval $i+1$ is
$\eta_{gd,i+1}$, not $\eta_{gd,i}$; this corrects the index typo
in the textbook's displayed (4.10).
Repeatedly moving all nominal rotations left gives
\begin{equation}
\Delta R_{ij}\simeq\Delta\tilde R_{ij}
\prod_{k=i}^{j-1}
\Exp\!\left(-\Delta\tilde R_{k+1,j}^{T}
J_{r,k}\eta_{gd,k}\Delta t_k\right).
\end{equation}
Combining the small noise exponentials to first order, define
\begin{equation}\label{eq:dphi}
\boxed{\delta\phi_{ij}\defeq
\sum_{k=i}^{j-1}\Delta\tilde R_{k+1,j}^{T}
J_{r,k}\eta_{gd,k}\Delta t_k,}
\end{equation}
where $\Delta\tilde R_{jj}=I_3$. Consequently,
\begin{equation}\label{eq:rotNoise}
\boxed{\Delta R_{ij}\simeq
\Delta\tilde R_{ij}\Exp(-\delta\phi_{ij}).}
\end{equation}
The adjoint rearrangement \eqref{eq:adjmove} is exact; the approximation
enters when separating and combining the noise exponentials.

\subsection{Velocity measurement and accumulated noise}
The relative velocity increment, with fixed accelerometer bias, is
\begin{equation}
\Delta v_{ij}
=\sum_{k=i}^{j-1}\Delta R_{ik}((\tilde a_k-b_{a,i})-\eta_{ad,k})\Delta t_k.
\end{equation}
Using $\Delta R_{ik}\simeq\Delta\tilde R_{ik}\Exp(-\delta\phi_{ik})$
and $\Exp(-\delta\phi_{ik})\simeq I-\delta\phi_{ik}^{\wedge}$,
\begin{align}
\Delta v_{ij}
&\simeq\sum_{k=i}^{j-1}
\Delta\tilde R_{ik}(I-\delta\phi_{ik}^{\wedge})
((\tilde a_k-b_{a,i})-\eta_{ad,k})\Delta t_k\notag\\
&\simeq\sum_{k=i}^{j-1}
\left[
\Delta\tilde R_{ik}(\tilde a_k-b_{a,i})
+\Delta\tilde R_{ik}(\tilde a_k-b_{a,i})^\wedge\delta\phi_{ik}
-\Delta\tilde R_{ik}\eta_{ad,k}
\right]\Delta t_k.
\end{align}
The sign of the rotation-noise term is \emph{positive} here, since
$-\delta\phi^\wedge f=f^\wedge\delta\phi$.
Define
\begin{equation}\label{eq:tildev}
\boxed{\Delta\tilde v_{ij}\defeq
\sum_{k=i}^{j-1}\Delta\tilde R_{ik}(\tilde a_k-b_{a,i})\Delta t_k,}
\end{equation}
and
\begin{equation}\label{eq:dvnoise}
\boxed{\delta v_{ij}\defeq
\sum_{k=i}^{j-1}\left[
-\Delta\tilde R_{ik}(\tilde a_k-b_{a,i})^\wedge\delta\phi_{ik}
+\Delta\tilde R_{ik}\eta_{ad,k}
\right]\Delta t_k.}
\end{equation}
Then
\begin{equation}\label{eq:velNoise}
\boxed{\Delta v_{ij}\simeq\Delta\tilde v_{ij}-\delta v_{ij}.}
\end{equation}
This sign convention agrees with the textbook's (4.14):
the cross-product order determines the apparent sign.

\subsection{Position measurement and accumulated noise}
Start with the previously derived position increment:
\begin{equation}
\Delta p_{ij}=
\sum_{k=i}^{j-1}
\left[
\Delta v_{ik}\Delta t_k+
\tfrac12\Delta R_{ik}((\tilde a_k-b_{a,i})-\eta_{ad,k})\Delta t_k^2
\right].
\end{equation}
Substitute $\Delta v_{ik}\simeq\Delta\tilde v_{ik}-\delta v_{ik}$
and $\Delta R_{ik}\simeq\Delta\tilde R_{ik}(I-\delta\phi_{ik}^\wedge)$:
\begin{align}
\Delta p_{ij}\simeq\sum_{k=i}^{j-1}\Big[
&\Delta\tilde v_{ik}\Delta t_k
+\tfrac12\Delta\tilde R_{ik}(\tilde a_k-b_{a,i})\Delta t_k^2
-\delta v_{ik}\Delta t_k\notag\\
&+\tfrac12\Delta\tilde R_{ik}(\tilde a_k-b_{a,i})^\wedge
\delta\phi_{ik}\Delta t_k^2
-\tfrac12\Delta\tilde R_{ik}\eta_{ad,k}\Delta t_k^2
\Big].
\end{align}
Define the computed pre-integrated position measurement
\begin{equation}\label{eq:tildep}
\boxed{\Delta\tilde p_{ij}\defeq
\sum_{k=i}^{j-1}\left[
\Delta\tilde v_{ik}\Delta t_k+
\tfrac12\Delta\tilde R_{ik}(\tilde a_k-b_{a,i})\Delta t_k^2
\right],}
\end{equation}
and its accumulated noise
\begin{equation}\label{eq:dpnoise}
\boxed{\begin{aligned}
\delta p_{ij}\defeq\sum_{k=i}^{j-1}\Big[
&\delta v_{ik}\Delta t_k
-\tfrac12\Delta\tilde R_{ik}(\tilde a_k-b_{a,i})^\wedge
\delta\phi_{ik}\Delta t_k^2\\
&+\tfrac12\Delta\tilde R_{ik}\eta_{ad,k}\Delta t_k^2
\Big].
\end{aligned}}
\end{equation}
Hence
\begin{equation}\label{eq:posNoise}
\boxed{\Delta p_{ij}\simeq\Delta\tilde p_{ij}-\delta p_{ij}.}
\end{equation}

\section{Summary: pre-integration measurement model}
Combining the preceding definitions with the relative-state increments
yields the following pre-integration measurement model:
\begin{subequations}\label{eq:measurementSummary}
\begin{align}
\boxed{\Delta\tilde R_{ij}}
&\simeq R_i^TR_j\Exp(\delta\phi_{ij}),\\
\boxed{\Delta\tilde v_{ij}}
&\simeq R_i^T(v_j-v_i-g\Delta t_{ij})+\delta v_{ij},\\
\boxed{\Delta\tilde p_{ij}}
&\simeq R_i^T\left(p_j-p_i-v_i\Delta t_{ij}
-\tfrac12g\Delta t_{ij}^2\right)+\delta p_{ij}.
\end{align}
\end{subequations}
The positive sign in the rotation measurement is consistent with
$\Delta R_{ij}\simeq\Delta\tilde R_{ij}\Exp(-\delta\phi_{ij})$:
right-multiply by $\Exp(\delta\phi_{ij})$ to isolate
$\Delta\tilde R_{ij}$. Unlike $\delta v_{ij}$ and $\delta p_{ij}$,
rotation noise enters multiplicatively on $\SO(3)$.

The three tilde quantities are computed from IMU measurements and the
fixed bias linearization point. The three delta noise quantities represent
accumulated measurement noise, whose covariance is propagated in the
following section. The bias-update Jacobians are derived subsequently.


\section{Recursive propagation of pre-integration noise}
\subsection{Conventions and initialization}
The accumulated noise is ordered as
\begin{equation}
\eta_{ik}\defeq\begin{bmatrix}\delta\phi_{ik}\\\delta v_{ik}\\\delta p_{ik}\end{bmatrix}\in\mathbb R^9,
\qquad
\eta_{d,k}\defeq\begin{bmatrix}\eta_{gd,k}\\\eta_{ad,k}\end{bmatrix}\in\mathbb R^6.
\end{equation}
We use the fixed biases $b_{g,i},b_{a,i}$ from the preceding section, and the sample interval $\Delta t_k=t_{k+1}-t_k$. The only additional shorthand is
\begin{equation}
J_{r,k}\defeq J_r\!\left((\tilde\omega_k-b_{g,i})\Delta t_k\right).
\end{equation}
The initial accumulated errors and covariance vanish:
\begin{equation}
\eta_{ii}=0,\qquad \Sigma_{ii}=0_{9\times9}.
\end{equation}

\subsection{Rotation noise recursion}
The accumulated rotation perturbation from the preceding section, with the
same general interval indices $i,j$ as in the textbook, is
\begin{equation}
\delta\phi_{ij}
=\sum_{\ell=i}^{j-1} \Delta\tilde R_{\ell+1,j}^{T}
 J_{r,\ell}\eta_{gd,\ell}\Delta t_\ell.
\end{equation}
To derive the recursive update, set $j=k+1$ in this general expression.
Separate the final term and use
$\Delta\tilde R_{\ell+1,k+1}
=\Delta\tilde R_{\ell+1,k}\Delta\tilde R_{k,k+1}$ for $\ell<k$:
\begin{align}
\delta\phi_{i,k+1}
&=\sum_{\ell=i}^{k-1}\Delta\tilde R_{k,k+1}^T\Delta\tilde R_{\ell+1,k}^{T}
 J_{r,\ell}\eta_{gd,\ell}\Delta t_\ell
 +J_{r,k}\eta_{gd,k}\Delta t_k\\
&=\boxed{\Delta\tilde R_{k,k+1}^T\delta\phi_{ik}+J_{r,k}\eta_{gd,k}\Delta t_k}.
\end{align}
The transpose of $\Delta\tilde R_{k,k+1}$ is the adjoint transport of the previously
accumulated rotation noise to the newest right-tangent coordinates.

\subsection{Velocity noise recursion}
For the general pre-integration interval $[t_i,t_j]$, the accumulated
velocity noise is
\begin{equation}
\delta v_{ij}
=\sum_{\ell=i}^{j-1}\left[
-\Delta\tilde R_{i\ell}(\tilde a_\ell-b_{a,i})^\wedge
\delta\phi_{i\ell}\Delta t_\ell
+\Delta\tilde R_{i\ell}\eta_{ad,\ell}\Delta t_\ell
\right].
\end{equation}
To obtain a recursive update, set $j=k+1$ and separate the final
IMU interval:
\begin{align}
\delta v_{i,k+1}
&=\sum_{\ell=i}^{k-1}\left[
-\Delta\tilde R_{i\ell}(\tilde a_\ell-b_{a,i})^\wedge
\delta\phi_{i\ell}
+\Delta\tilde R_{i\ell}\eta_{ad,\ell}\right]\Delta t_\ell\notag\\
&\quad-\Delta\tilde R_{ik}(\tilde a_k-b_{a,i})^\wedge
\delta\phi_{ik}\Delta t_k
+\Delta\tilde R_{ik}\eta_{ad,k}\Delta t_k\\
&=\delta v_{ik}
-\Delta\tilde R_{ik}(\tilde a_k-b_{a,i})^\wedge
\delta\phi_{ik}\Delta t_k
+\Delta\tilde R_{ik}\eta_{ad,k}\Delta t_k.
\end{align}
The velocity error therefore couples to the accumulated rotation error.

\subsection{Position noise recursion}
For the general interval $[t_i,t_j]$, the accumulated position noise is
\begin{equation}
\delta p_{ij}=\sum_{\ell=i}^{j-1}\bigl[\delta v_{i\ell}\Delta t_\ell-\tfrac12\Delta\tilde R_{i\ell}(\tilde a_\ell-b_{a,i})^\wedge\delta\phi_{i\ell}\Delta t_\ell^2+\tfrac12\Delta\tilde R_{i\ell}\eta_{ad,\ell}\Delta t_\ell^2\bigr].
\end{equation}
Setting $j=k+1$ and splitting off the final interval yields
\begin{equation}
\delta p_{i,k+1}=\delta p_{ik}+\delta v_{ik}\Delta t_k-\tfrac12\Delta\tilde R_{ik}(\tilde a_k-b_{a,i})^\wedge\delta\phi_{ik}\Delta t_k^2+\tfrac12\Delta\tilde R_{ik}\eta_{ad,k}\Delta t_k^2.
\end{equation}
Notice that the $\delta v_{ik}$ term is the \emph{previous} accumulated
velocity noise, consistent with the discrete position integration rule.

\subsection{Nine-dimensional matrix recursion}
Collecting the preceding three relations gives
\begin{equation}\label{eq:noise_recursion}
\boxed{\eta_{i,k+1}=A_k\eta_{ik}+B_k\eta_{d,k}},
\end{equation}
where the $9\times9$ transition matrix is
\begin{equation}\label{eq:A_noise}
\boxed{
A_k=\begin{bmatrix}
\Delta\tilde R_{k,k+1}^T&0&0\\
-\Delta\tilde R_{ik} (\tilde a_k-b_{a,i})^\wedge \Delta t_k&I_3&0\\
-\tfrac12 \Delta\tilde R_{ik} (\tilde a_k-b_{a,i})^\wedge \Delta t_k^2&I_3\Delta t_k&I_3
\end{bmatrix}}
\end{equation}
and the $9\times6$ noise-input matrix is
\begin{equation}\label{eq:B_noise}
\boxed{
B_k=\begin{bmatrix}
J_{r,k}\Delta t_k&0\\
0&\Delta\tilde R_{ik}\Delta t_k\\
0&\tfrac12 \Delta\tilde R_{ik}\Delta t_k^2
\end{bmatrix}.}
\end{equation}
Rows are ordered $(\delta\phi,\delta v,\delta p)$ and columns of
$B_k$ are ordered $(\eta_{gd},\eta_{ad})$. These blocks are the
coefficients of the preceding three recursions, not independent
assumptions. We have used the interval index $k$ consistently; the
book sometimes labels the last interval $j-1$.

\section{Covariance propagation of pre-integrated measurements}
\subsection{Discrete IMU noise covariance}
Let $Q_{d,k}\defeq\operatorname{Cov}(\eta_{d,k})$.
If gyroscope and accelerometer discrete noise are independent,
\begin{equation}
Q_{d,k}=\begin{bmatrix}
Q_{g,d,k}&0\\0&Q_{a,d,k}
\end{bmatrix}.
\end{equation}
With continuous-time white-noise covariance densities $Q_{g,c}$ and
$Q_{a,c}$, the interval-average discrete noise convention used in the
textbook gives
\begin{equation}
Q_{g,d,k}=\frac{Q_{g,c}}{\Delta t_k},\qquad
Q_{a,d,k}=\frac{Q_{a,c}}{\Delta t_k}.
\end{equation}
Thus the covariance of the \emph{noise increments} entering the
rotation/velocity equations scales with $\Delta t_k$ after multiplication
by $B_k$. Do not confuse this $Q_{d,k}$ with the accumulated
pre-integration covariance $\Sigma_{ik}$.

\subsection{Derivation from the recursion}
Define $\Sigma_{ik}=\operatorname{Cov}(\eta_{ik})$.
Assume the new IMU noise $\eta_{d,k}$ is zero-mean and uncorrelated
with the previously accumulated error $\eta_{ik}$; treat $A_k,B_k$
as fixed at the current linearization point. Then
\begin{align}
\Sigma_{i,k+1}
&=\operatorname{Cov}(A_k\eta_{ik}+B_k\eta_{d,k})\\
&=A_k\Sigma_{ik}A_k^T+B_kQ_{d,k}B_k^T,
\end{align}
because the two cross-covariance terms vanish. Therefore,
\begin{equation}\label{eq:preint_cov}
\boxed{\Sigma_{i,k+1}=A_k\Sigma_{ik}A_k^T+B_kQ_{d,k}B_k^T,
\qquad \Sigma_{ii}=0.}
\end{equation}

\subsection{Pre-integration propagation flow}
\[
\boxed{\begin{array}{c}
\text{Receive IMU sample }(\tilde\omega_k,\tilde a_k,\Delta t_k)\\[2mm]
\downarrow\\[2mm]
\text{Update }\Delta\tilde R_{i,k+1},\Delta\tilde v_{i,k+1},\Delta\tilde p_{i,k+1}\\[2mm]
\downarrow\\[2mm]
\text{Compute }A_k,\ B_k,\ Q_{d,k}\\[2mm]
\downarrow\\[2mm]
\Sigma_{i,k+1}=A_k\Sigma_{ik}A_k^T+B_kQ_{d,k}B_k^T\\[2mm]
\downarrow\\[2mm]
\text{Continue until endpoint }j
\end{array}}
\]
At the endpoint, the computed measurements
$(\Delta\tilde R_{ij},\Delta\tilde v_{ij},\Delta\tilde p_{ij})$
and their $9\times9$ covariance $\Sigma_{ij}$ describe the
pre-integrated IMU measurement and its uncertainty. This completes the measurement-noise and covariance propagation of
the textbook's Section 4.2.2. The next section begins the bias-update model.


\section{Bias Update: First-Order Measurement Correction}
\subsection{General first-order approximation of a rotation-valued function}
Let $R:\R^3\to\SO(3)$ be a differentiable rotation-valued function
of a parameter $\phi\in\R^3$. For a sufficiently small parameter
perturbation $\delta\phi$, define its \emph{right-local rotation increment}
$f(\delta\phi)$ by
\begin{equation}\label{eq:rot_param_def}
\boxed{R(\phi+\delta\phi)=R(\phi)\Exp\bigl(f(\delta\phi)\bigr).}
\end{equation}
This is an exact local parameterization of the relative rotation, not
yet a first-order approximation. Multiplying on the left by $R(\phi)^T$
and applying the local logarithm gives
\begin{equation}\label{eq:rot_param_f}
\boxed{f(\delta\phi)=\Log\!\left(R(\phi)^T R(\phi+\delta\phi)\right).}
\end{equation}
Because $f(0)=\Log(I)=0$, its ordinary Euclidean first-order Taylor
expansion at the zero perturbation is
\begin{equation}
\begin{aligned}
f(\delta\phi)
&=f(0)+\left.\frac{\partial f}{\partial\delta\phi}\right|_{\delta\phi=0}
\delta\phi+O(\|\delta\phi\|^2)\\
&=J_{R,\phi}\,\delta\phi+O(\|\delta\phi\|^2),
\end{aligned}
\end{equation}
where the \emph{right-trivialized parameter Jacobian} is defined by
\begin{equation}\label{eq:rot_param_J}
\boxed{J_{R,\phi}\defeq
\left.\frac{\partial}{\partial\delta\phi}
\Log\!\left(R(\phi)^T R(\phi+\delta\phi)\right)
\right|_{\delta\phi=0}.}
\end{equation}
Substituting into \eqref{eq:rot_param_def} yields the first-order
rotation-valued approximation
\begin{equation}\label{eq:rot_param_linear}
\boxed{R(\phi+\delta\phi)\simeq
R(\phi)\Exp(J_{R,\phi}\delta\phi).}
\end{equation}
Here $\phi$ is a generic \emph{input parameter}; it is not the
pre-integration rotation-noise variable $\delta\phi_{ij}$.

\subsection{Application to the pre-integrated rotation measurement}
In the textbook's bias-update equation, set the generic input parameter
to the initial gyroscope-bias estimate $b_{g,i}$. The pre-integrated
rotation $\Delta\tilde R_{ij}$ is a function of this bias because
its IMU increments contain $\tilde\omega_k-b_{g,i}$.
Thus, applying the right-trivialized Jacobian defined in
\eqref{eq:rot_param_J} at the fixed bias $b_{g,i}$ gives
\begin{equation}\label{eq:preint_bg_rotation}
\boxed{\Delta\tilde R_{ij}(b_{g,i}+\delta b_{g,i})\simeq
\Delta\tilde R_{ij}(b_{g,i})
\Exp\!\left(J_{\Delta\tilde R,b_g}^{ij}\delta b_{g,i}\right).}
\end{equation}
Here $J_{\Delta\tilde R,b_g}^{ij}$ denotes $J_{R,\phi}$ from
\eqref{eq:rot_param_J} evaluated for
$R(\phi)=\Delta\tilde R_{ij}(b_{g,i})$ with $\phi=b_{g,i}$.
Equivalently, it is the right-trivialized $3\times3$ bias Jacobian
that the textbook denotes by
$\partial\Delta\tilde R_{ij}/\partial b_{g,i}$; it is not an
entrywise derivative of the rotation matrix. The subsequent
bias-Jacobian derivations retain the textbook notation for comparison.

\subsection{Velocity and position bias corrections}
Unlike the rotation measurement, $\Delta\tilde v_{ij}$ and
$\Delta\tilde p_{ij}$ are vectors in $\R^3$, so ordinary first-order
Taylor expansions apply. Let the fixed bias linearization point be
$(b_{g,i},b_{a,i})$ and its correction be
$(\delta b_{g,i},\delta b_{a,i})$. Then
\begin{equation}\label{eq:preint_bg_velocity}
\boxed{\begin{aligned}
&\Delta\tilde v_{ij}(b_{g,i}+\delta b_{g,i},
 b_{a,i}+\delta b_{a,i})\\
&\quad\simeq\Delta\tilde v_{ij}(b_{g,i},b_{a,i})
+\frac{\partial\Delta\tilde v_{ij}}{\partial b_{g,i}}\delta b_{g,i}
+\frac{\partial\Delta\tilde v_{ij}}{\partial b_{a,i}}\delta b_{a,i}.
\end{aligned}}
\end{equation}
Likewise,
\begin{equation}\label{eq:preint_bg_position}
\boxed{\begin{aligned}
&\Delta\tilde p_{ij}(b_{g,i}+\delta b_{g,i},
 b_{a,i}+\delta b_{a,i})\\
&\quad\simeq\Delta\tilde p_{ij}(b_{g,i},b_{a,i})
+\frac{\partial\Delta\tilde p_{ij}}{\partial b_{g,i}}\delta b_{g,i}
+\frac{\partial\Delta\tilde p_{ij}}{\partial b_{a,i}}\delta b_{a,i}.
\end{aligned}}
\end{equation}
All five bias Jacobians are evaluated at the fixed bias linearization
point. The rotation bias Jacobian $J_{\Delta\tilde R,b_g}^{ij}$,
defined using the right-trivialized derivative in
\eqref{eq:rot_param_J}, maps a gyroscope bias perturbation to a local
rotation vector in $\R^3$. The remaining four Jacobians, associated
with the pre-integrated velocity and position, are ordinary Euclidean
derivatives.

Equations \eqref{eq:preint_bg_rotation}--\eqref{eq:preint_bg_position}
summarize the three first-order bias-correction relationships in the
textbook's (4.27). The rotation correction uses the exponential map to
remain on $\SO(3)$, whereas the velocity and position corrections use
ordinary vector addition. The five bias Jacobians are derived in the
following sections.

\subsection{Gyroscope-bias Jacobian: product expansion}
We now derive the rotation correction corresponding to the textbook's
(4.28)--(4.29), using the same variable sampling intervals $\Delta t_k$
as in the preceding sections. Throughout the pre-integration interval,
$b_{g,i}$ is the fixed bias linearization point, and $\delta b_{g,i}$
is a small change to that \emph{same} bias for every IMU sample.

Start from the pre-integrated rotation as a function of the bias:
\begin{equation}
\Delta\tilde R_{ij}(b_{g,i}+\delta b_{g,i})
=\prod_{k=i}^{j-1}\Exp\!\left(
(\tilde\omega_k-b_{g,i}-\delta b_{g,i})\Delta t_k\right).
\end{equation}
For each sample, the right-Jacobian increment formula
from BCH yields
\begin{equation}\label{eq:bias_one_step}
\Exp((\tilde\omega_k-b_{g,i}-\delta b_{g,i})\Delta t_k)
\simeq\Exp((\tilde\omega_k-b_{g,i})\Delta t_k)\Exp(-J_{r,k}\delta b_{g,i}\Delta t_k).
\end{equation}
The Jacobian $J_{r,k}$ must be retained because $(\tilde\omega_k-b_{g,i})\Delta t_k$ is the
nominal increment, not a perturbation that vanishes with $\delta b_{g,i}$.

As in the earlier rotation-noise derivation, use the exact adjoint
identity $\Exp(\epsilon)B=B\Exp(B^T\epsilon)$ to move each small
bias perturbation to the right. First-order BCH then combines these
perturbations by addition, neglecting their products, which are
$O(\|\delta b_{g,i}\|^2)$. Thus the perturbation introduced at sample
$k$ is transported by the subsequent nominal rotation
$\Delta\tilde R_{k+1,j}^{T}$.

Repeating the same adjoint transport across the remaining increments,
with $\Delta\tilde R_{jj}=I_3$, gives
\begin{equation}\label{eq:bias_rotation_product}
\boxed{\begin{aligned}
&\Delta\tilde R_{ij}(b_{g,i}+\delta b_{g,i})\\
&\quad\simeq \Delta\tilde R_{ij}(b_{g,i})
\Exp\!\left[-\left(\sum_{k=i}^{j-1}
\Delta\tilde R_{k+1,j}^{T}J_{r,k}\Delta t_k\right)
\delta b_{g,i}\right].
\end{aligned}}
\end{equation}
This is the first-order rotation bias correction in the textbook's (4.28).

Comparing \eqref{eq:bias_rotation_product} with the previously defined
right-trivialized linearization
$\Delta\tilde R_{ij}(b+\delta b)\simeq
\Delta\tilde R_{ij}(b)\Exp(\frac{\partial\Delta\tilde R_{ij}}{\partial b_{g,i}}\delta b)$ gives
\begin{equation}\label{eq:bias_rotation_jac_sum}
\boxed{\frac{\partial\Delta\tilde R_{ij}}{\partial b_{g,i}}
=-\sum_{k=i}^{j-1}\Delta\tilde R_{k+1,j}^{T}
J_{r,k}\Delta t_k.}
\end{equation}
The derivative notation above refers to the \emph{right-trivialized}
$3\times3$ Jacobian, not an elementwise derivative of a rotation matrix.

\subsection{Recursive form of the gyroscope-bias Jacobian}
To obtain the textbook's (4.30), separate the last term of
\eqref{eq:bias_rotation_jac_sum}:
\begin{align}
\frac{\partial\Delta\tilde R_{ij}}{\partial b_{g,i}}
&=-\sum_{k=i}^{j-2}\Delta\tilde R_{k+1,j}^{T}
J_{r,k}\Delta t_k
-\Delta\tilde R_{j,j}^{T}J_{r,j-1}\Delta t_{j-1}\notag\\
&=-\sum_{k=i}^{j-2}\Delta\tilde R_{k+1,j}^{T}
J_{r,k}\Delta t_k-J_{r,j-1}\Delta t_{j-1}.
\end{align}
For $k\le j-2$, the nominal relative rotations compose as
\begin{equation}
\Delta\tilde R_{k+1,j}
=\Delta\tilde R_{k+1,j-1}\Delta\tilde R_{j-1,j},
\qquad
\Delta\tilde R_{k+1,j}^{T}
=\Delta\tilde R_{j-1,j}^{T}\Delta\tilde R_{k+1,j-1}^{T}.
\end{equation}
Substituting and factoring out the final rotation gives
\begin{align}
\frac{\partial\Delta\tilde R_{ij}}{\partial b_{g,i}}
&=\Delta\tilde R_{j-1,j}^{T}
\left[-\sum_{k=i}^{j-2}\Delta\tilde R_{k+1,j-1}^{T}
J_{r,k}\Delta t_k\right]-J_{r,j-1}\Delta t_{j-1}\notag\\
&=\boxed{\Delta\tilde R_{j-1,j}^{T}
\frac{\partial\Delta\tilde R_{i,j-1}}{\partial b_{g,i}}-J_{r,j-1}\Delta t_{j-1}}.
\end{align}
Equivalently, for an incoming IMU interval $[t_k,t_{k+1}]$,
\begin{equation}\label{eq:bias_rotation_jac_rec}
\boxed{\frac{\partial\Delta\tilde R_{i,k+1}}{\partial b_{g,i}}
=\Delta\tilde R_{k,k+1}^{T}\frac{\partial\Delta\tilde R_{ik}}{\partial b_{g,i}}
-J_{r,k}\Delta t_k,\qquad \frac{\partial\Delta\tilde R_{ii}}{\partial b_{g,i}}=0_{3\times3}.}
\end{equation}
The first term transports the previously accumulated Jacobian to the
new rotation reference, and the second adds the current interval's
bias sensitivity. In the textbook's (4.30), the last $J_{r,k}\Delta t$
should be read as the contribution of the \emph{last} interval,
$J_{r,j-1}\Delta t_{j-1}$, when the left-hand side ends at $j$.


\section{Velocity and Position Bias Jacobians}
We now complete the five bias Jacobians of the pre-integrated
measurements, corresponding to the textbook's (4.33)--(4.34).
All Jacobians are evaluated at the fixed bias linearization point
$(b_{g,i},b_{a,i})$. We retain the sampling interval $\Delta t_k$.

\subsection{Velocity with respect to accelerometer bias}
From the measurement definition,
\begin{equation}
\Delta\tilde v_{ij}=\sum_{k=i}^{j-1}
\Delta\tilde R_{ik}(\tilde a_k-b_{a,i})\Delta t_k.
\end{equation}
The rotation increment does not depend on accelerometer bias. Thus
\begin{equation}\label{eq:Jvba_sum}
\boxed{\frac{\partial\Delta\tilde v_{ij}}{\partial b_{a,i}}
=-\sum_{k=i}^{j-1}\Delta\tilde R_{ik}\Delta t_k.}
\end{equation}

\subsection{Velocity with respect to gyroscope bias}
The gyro bias affects the velocity through the pre-integrated
rotation. Apply the right-trivialized perturbation derived earlier:
\begin{equation}
\Delta\tilde R_{ik}(b_{g,i}+\delta b_g)
\simeq\Delta\tilde R_{ik}
\Exp(\frac{\partial\Delta\tilde R_{ik}}{\partial b_{g,i}}\delta b_g).
\end{equation}
Keeping only first-order terms,
\begin{align}
\Delta\tilde v_{ij}(b_{g,i}+\delta b_g)
&\simeq\sum_{k=i}^{j-1}
\Delta\tilde R_{ik}
\left[I+(\frac{\partial\Delta\tilde R_{ik}}{\partial b_{g,i}}\delta b_g)^\wedge\right]
(\tilde a_k-b_{a,i})\Delta t_k\notag\\
&=\Delta\tilde v_{ij}(b_{g,i})
-\sum_{k=i}^{j-1}
\Delta\tilde R_{ik}(\tilde a_k-b_{a,i})^\wedge
\frac{\partial\Delta\tilde R_{ik}}{\partial b_{g,i}}\delta b_g\Delta t_k.
\end{align}
Here $u^\wedge w=-w^\wedge u$. Comparing with the first-order
Taylor expansion gives
\begin{equation}\label{eq:Jvbg_sum}
\boxed{\frac{\partial\Delta\tilde v_{ij}}{\partial b_{g,i}}
=-\sum_{k=i}^{j-1}
\Delta\tilde R_{ik}(\tilde a_k-b_{a,i})^\wedge
\frac{\partial\Delta\tilde R_{ik}}{\partial b_{g,i}}\Delta t_k.}
\end{equation}

\subsection{Position with respect to accelerometer bias}
Recall
\begin{equation}
\Delta\tilde p_{ij}
=\sum_{k=i}^{j-1}\left[
\Delta\tilde v_{ik}\Delta t_k+
\frac12\Delta\tilde R_{ik}
(\tilde a_k-b_{a,i})\Delta t_k^2\right].
\end{equation}
Differentiate both terms. The rotation has no dependence on $b_{a,i}$:
\begin{equation}\label{eq:Jpba_sum}
\boxed{\frac{\partial\Delta\tilde p_{ij}}{\partial b_{a,i}}
=\sum_{k=i}^{j-1}
\left[\frac{\partial\Delta\tilde v_{ik}}{\partial b_{a,i}}\Delta t_k
-\frac12\Delta\tilde R_{ik}\Delta t_k^2\right].}
\end{equation}

\subsection{Position with respect to gyroscope bias}
There are two dependency paths: gyro bias affects the accumulated
velocity $\Delta\tilde v_{ik}$ and the rotation
$\Delta\tilde R_{ik}$ in the acceleration contribution.
Under $b_{g,i}\mapsto b_{g,i}+\delta b_g$, insert
\begin{align}
\Delta\tilde v_{ik}(b_{g,i}+\delta b_g)
&\simeq\Delta\tilde v_{ik}+\frac{\partial\Delta\tilde v_{ik}}{\partial b_{g,i}}\delta b_g,\\
\Delta\tilde R_{ik}(b_{g,i}+\delta b_g)
&\simeq\Delta\tilde R_{ik}
\left[I+(\frac{\partial\Delta\tilde R_{ik}}{\partial b_{g,i}}\delta b_g)^\wedge\right].
\end{align}
After collecting the terms linear in $\delta b_g$,
\begin{align}
\Delta\tilde p_{ij}(b_{g,i}+\delta b_g)
&\simeq\Delta\tilde p_{ij}(b_{g,i})\notag\\
&\quad+\sum_{k=i}^{j-1}\left[
\frac{\partial\Delta\tilde v_{ik}}{\partial b_{g,i}}\Delta t_k
-\frac12\Delta\tilde R_{ik}
(\tilde a_k-b_{a,i})^\wedge
\frac{\partial\Delta\tilde R_{ik}}{\partial b_{g,i}}\Delta t_k^2
\right]\delta b_g.
\end{align}
Therefore,
\begin{equation}\label{eq:Jpbg_sum}
\begingroup\small
\boxed{\displaystyle
\frac{\partial\Delta\tilde p_{ij}}{\partial b_{g,i}}
=\sum_{k=i}^{j-1}\!\left[
\frac{\partial\Delta\tilde v_{ik}}{\partial b_{g,i}}\Delta t_k
-\frac12\Delta\tilde R_{ik}(\tilde a_k-b_{a,i})^\wedge
\frac{\partial\Delta\tilde R_{ik}}{\partial b_{g,i}}\Delta t_k^2
\right].}
\endgroup
\end{equation}

\section{Recursive Forms of the Five Bias Jacobians}
For a new IMU interval $[t_k,t_{k+1}]$, split the last term
from each sum. The rotation Jacobian has already been derived.
The complete set of recursions is
\begin{align}
\frac{\partial\Delta\tilde R_{i,k+1}}{\partial b_{g,i}}
&=\Delta\tilde R_{k,k+1}^{T}\frac{\partial\Delta\tilde R_{ik}}{\partial b_{g,i}}
-J_{r,k}\Delta t_k,\label{eq:recRbg}\\
\frac{\partial\Delta\tilde v_{i,k+1}}{\partial b_{a,i}}
&=\frac{\partial\Delta\tilde v_{ik}}{\partial b_{a,i}}-\Delta\tilde R_{ik}\Delta t_k,
\label{eq:recVba}\\
\frac{\partial\Delta\tilde v_{i,k+1}}{\partial b_{g,i}}
&=\frac{\partial\Delta\tilde v_{ik}}{\partial b_{g,i}}
-\Delta\tilde R_{ik}(\tilde a_k-b_{a,i})^\wedge
\frac{\partial\Delta\tilde R_{ik}}{\partial b_{g,i}}\Delta t_k,\label{eq:recVbg}\\
\frac{\partial\Delta\tilde p_{i,k+1}}{\partial b_{a,i}}
&=\frac{\partial\Delta\tilde p_{ik}}{\partial b_{a,i}}+\frac{\partial\Delta\tilde v_{ik}}{\partial b_{a,i}}\Delta t_k
-\frac12\Delta\tilde R_{ik}\Delta t_k^2,
\label{eq:recPba}\\
\frac{\partial\Delta\tilde p_{i,k+1}}{\partial b_{g,i}}
&=\frac{\partial\Delta\tilde p_{ik}}{\partial b_{g,i}}+\frac{\partial\Delta\tilde v_{ik}}{\partial b_{g,i}}\Delta t_k\notag\\
&\quad-\frac12\Delta\tilde R_{ik}
(\tilde a_k-b_{a,i})^\wedge
\frac{\partial\Delta\tilde R_{ik}}{\partial b_{g,i}}\Delta t_k^2.
\label{eq:recPbg}
\end{align}
All five Jacobians are initially zero at $k=i$.
The right-hand sides use the \emph{old} values at $k$,
including $\Delta\tilde R_{ik}$ and the old velocity and rotation
Jacobians. In an in-place implementation, update the position
Jacobians first, then the velocity Jacobians, and finally the
rotation Jacobian, or cache all old quantities.

\subsection{Summary}
The rotation bias Jacobian is a right-trivialized $\SO(3)$
derivative, whereas the other four bias Jacobians are ordinary
Euclidean derivatives. Equations \eqref{eq:recRbg}--\eqref{eq:recPbg}
provide recursive updates for the sensitivities of the pre-integrated
rotation, velocity, and position measurements with respect to the
gyroscope and accelerometer biases.

These Jacobians enable efficient first-order bias correction without
re-integrating the original IMU measurements whenever the bias estimates
change during graph optimization. Together with the pre-integrated
measurements and their covariance, they provide the information needed
to construct the IMU pre-integration factor in the following section.

\section{Pre-integration Factor in Graph Optimization}
\subsection{Purpose and keyframe states}
The pre-integration measurement model derived above can now be used as
an IMU constraint between keyframes $i$ and $j$.  At each keyframe,
we use the state variables
\begin{equation}
 x_k=(R_k,p_k,v_k,b_{a,k},b_{g,k}).
\end{equation}
The purpose of this section is to define the \emph{IMU pre-integration
factor}, not yet to assemble a complete multi-sensor optimization
problem. The pre-integrated measurements are
$\Delta\tilde R_{ij}$, $\Delta\tilde v_{ij}$ and
$\Delta\tilde p_{ij}$, with covariance $\Sigma_{ij}$.

\subsection{Residual definition}
The measurement model gives a predicted relative motion from the
keyframe states. The residual compares this prediction against the
pre-integrated IMU measurement. Following the textbook's residual
convention, define
\begin{subequations}\label{eq:imu_factor_residual}
\begin{align}
 r_{\Delta R_{ij}}
 &=\Log\!\left(\Delta\tilde R_{ij}^{T}R_i^TR_j\right),\\
 r_{\Delta v_{ij}}
 &=R_i^T(v_j-v_i-g\Delta t_{ij})-\Delta\tilde v_{ij},\\
 r_{\Delta p_{ij}}
 &=R_i^T\!\left(p_j-p_i-v_i\Delta t_{ij}
 -\tfrac12g\Delta t_{ij}^{2}\right)-\Delta\tilde p_{ij}.
\end{align}
\end{subequations}
The rotation residual uses the $SO(3)$ logarithm instead of
subtracting rotation matrices. When the state-predicted relative
motion agrees with the pre-integrated measurements, all three
residuals vanish.

Stack the three components in a nine-dimensional vector:
\begin{equation}
\boxed{r_{ij}^{\mathrm{IMU}}=
\begin{bmatrix}
 r_{\Delta R_{ij}}\\ r_{\Delta v_{ij}}\\ r_{\Delta p_{ij}}
\end{bmatrix}\in\mathbb R^9.}
\end{equation}
The pre-integration covariance derived previously supplies the
weighting for a least-squares IMU factor, schematically
\begin{equation}
 E_{ij}^{\mathrm{IMU}}
 =\left(r_{ij}^{\mathrm{IMU}}\right)^T
 \Sigma_{ij}^{-1}r_{ij}^{\mathrm{IMU}}.
\end{equation}
This is the standard weighted least-squares interpretation of the
textbook's residual and covariance; the section primarily defines
the residuals and factor connectivity.


\subsection{States, measurements, and residuals}
These are three distinct mathematical objects. The keyframe states
$x_i,x_j$ are unknown optimization variables (15 degrees of freedom
per keyframe). The pre-integrated IMU measurements
$(\Delta\tilde R_{ij},\Delta\tilde v_{ij},\Delta\tilde p_{ij})$
are obtained from IMU samples at a fixed bias linearization point;
they have nine local degrees of freedom, but are \emph{not} navigation
states being optimized. The measurement residual
$r_{ij}^{\mathrm{IMU}}\in\mathbb R^9$ compares the motion predicted
from the current keyframe states with these measurements.

A factor-graph connection means that a residual \emph{depends on}
a state variable; it does not imply that the IMU directly measures
that variable. The dependencies, including first-order bias
correction, are
\begin{align*}
 r_{\Delta R_{ij}}&=r_{\Delta R_{ij}}(R_i,R_j,b_{g,i}),\\
 r_{\Delta v_{ij}}&=r_{\Delta v_{ij}}(R_i,v_i,v_j,b_{g,i},b_{a,i}),\\
 r_{\Delta p_{ij}}&=r_{\Delta p_{ij}}(R_i,p_i,p_j,v_i,b_{g,i},b_{a,i}).
\end{align*}
The endpoint biases $b_{a,j},b_{g,j}$ do not enter this nine-dimensional
pre-integration residual. Instead, separate bias random-walk factors
constrain their changes, for example
\begin{equation}
 r_{b_a}=b_{a,j}-b_{a,i},\qquad
 r_{b_g}=b_{g,j}-b_{g,i}.
\end{equation}
These two factors each have three residual components.

\subsection{Jacobian dependency matrix}
The rows below represent \emph{measurement residuals}, not
pre-integrated measurements; the columns represent local perturbations
of optimization states. A check mark indicates a generally nonzero
$3\times3$ Jacobian block. A dash denotes a zero block.
\begin{center}
\begingroup\setlength{\tabcolsep}{4pt}\renewcommand{\arraystretch}{1.25}
\resizebox{\textwidth}{!}{$
\begin{array}{c|cccccccccc}
 &\phi_i&\delta p_i&\delta v_i&\delta b_{a,i}&\delta b_{g,i}
 &\phi_j&\delta p_j&\delta v_j&\delta b_{a,j}&\delta b_{g,j}\\\hline
 r_{\Delta R_{ij}}&\checkmark&-&-&-&\checkmark&\checkmark&-&-&-&-\\
 r_{\Delta v_{ij}}&\checkmark&-&\checkmark&\checkmark&\checkmark&-&-&\checkmark&-&-\\
 r_{\Delta p_{ij}}&\checkmark&\checkmark&\checkmark&\checkmark&\checkmark&-&\checkmark&-&-&-
\end{array}$}
\endgroup
\end{center}
Here $R_i'=R_i\Exp(\phi_i)$ and $R_j'=R_j\Exp(\phi_j)$.
For the complete two-keyframe local perturbation vector
$\delta X_{ij}\in\mathbb R^{30}$, the residual Jacobian is
\begin{equation}
 J_{\mathrm{pre}}=\frac{\partial r_{ij}^{\mathrm{IMU}}}
 {\partial\delta X_{ij}}\in\mathbb R^{9\times30}.
\end{equation}
The final two bias columns are zero. Restricting to the eight
three-dimensional state blocks on which the pre-integration residual
depends gives a $9\times24$ Jacobian. Its assembled block form is
recorded in the Appendix; the main text develops the individual
Jacobian blocks directly from the residual definitions.

\subsection{Dependence on the IMU biases}
Although the residual expressions explicitly display only the
rotation, position and velocity states, their pre-integrated
measurement terms depend on the biases at the starting keyframe:
\begin{align}
 \Delta\tilde R_{ij}&=\Delta\tilde R_{ij}(b_{g,i}),\\
 \Delta\tilde v_{ij}&=\Delta\tilde v_{ij}(b_{g,i},b_{a,i}),\\
 \Delta\tilde p_{ij}&=\Delta\tilde p_{ij}(b_{g,i},b_{a,i}).
\end{align}
The five bias Jacobians derived in the preceding section provide
first-order corrections to these measurements when the optimization
updates the biases, without reintegrating all IMU samples.
The IMU pre-integration factor therefore connects
$(R_i,p_i,v_i,b_{a,i},b_{g,i})$ to $(R_j,p_j,v_j)$.
The bias random walks between the two keyframes are represented by
separate bias factors, as illustrated in the textbook's Figure 4-1.

\subsection{Scope: IMU constraints versus other sensor factors}
At this stage, we have formulated \emph{only} the IMU pre-integration
constraint and the associated bias random-walk model. We have not yet
introduced GNSS, wheel-encoder or other sensor observation factors
into the graph. The IMU factor constrains relative motion, rather
than providing an absolute position or attitude reference by itself.

Later, a full graph may combine IMU pre-integration, bias random walk,
GNSS, odometry and prior factors. The precise graph construction and
implementation are addressed in later sections of the chapter.

\subsection{Summary}
The conceptual progression is
\[
\boxed{\text{IMU samples}\;\longrightarrow\;
(\Delta\tilde R_{ij},\Delta\tilde v_{ij},\Delta\tilde p_{ij},\Sigma_{ij})
\;\longrightarrow\;r_{ij}^{\mathrm{IMU}}
\;\longrightarrow\;\text{IMU factor}.}
\]
The following sections compute the residual Jacobians with respect to
the graph's optimization variables.


\section{Jacobians of Preintegration Residuals}
\subsection{Rotation residual: right perturbations}
We now linearize the rotation residual of the IMU factor with respect to
its endpoint rotations. Following the same convention as the preceding
sections, perturb each rotation on the right:
\begin{equation}
R_i' = R_i\Exp(\phi_i),\qquad
R_j' = R_j\Exp(\phi_j),
\end{equation}
where $\phi_i,\phi_j\in\mathbb R^3$ are small optimization increments,
not the pre-integration rotation noise $\delta\phi_{ij}$.
Recall the residual
\begin{equation}\label{eq:resrot_jac_base}
r_{\Delta R_{ij}}=\Log\!\left(\Delta\tilde R_{ij}^{T}R_i^TR_j\right).
\end{equation}
The pre-integrated measurement is held fixed when differentiating with
respect to the endpoint rotations.

\subsection{Derivative with respect to the initial rotation}
Perturb $R_i$ while holding $R_j$ fixed:
\begin{align}
r_{\Delta R_{ij}}(R_i\Exp(\phi_i))
&=\Log\!\left(\Delta\tilde R_{ij}^{T}
 (R_i\Exp(\phi_i))^TR_j\right)\notag\\
&=\Log\!\left(\Delta\tilde R_{ij}^{T}
 \Exp(-\phi_i)R_i^TR_j\right).
\end{align}
The $\SO(3)$ adjoint identity
\begin{equation}
\Exp(\xi)Q=Q\Exp(Q^T\xi),\qquad Q\in\SO(3),
\end{equation}
with $Q=R_i^TR_j$, moves the small perturbation to the right:
\begin{align}
\Exp(-\phi_i)R_i^TR_j
&=R_i^TR_j\Exp(-R_j^TR_i\phi_i),\\
r_{\Delta R_{ij}}(R_i\Exp(\phi_i))
&=\Log\!\left(\Delta\tilde R_{ij}^{T}R_i^TR_j
 \Exp(-R_j^TR_i\phi_i)\right).
\end{align}
By \eqref{eq:resrot_jac_base},
$\Delta\tilde R_{ij}^{T}R_i^TR_j=\Exp(r_{\Delta R_{ij}})$.
Use the first-order right-BCH identity
\begin{equation}\label{eq:right_bch_residual}
\Log\!\left(\Exp(a)\Exp(\epsilon)\right)
\simeq a+J_r^{-1}(a)\epsilon,
\qquad \epsilon\text{ small}.
\end{equation}
Consequently,
\begin{align}
r_{\Delta R_{ij}}(R_i\Exp(\phi_i))
&\simeq\Log\!\left(\Exp(r_{\Delta R_{ij}})
 \Exp(-R_j^TR_i\phi_i)\right)\notag\\
&\simeq r_{\Delta R_{ij}}
 -J_r^{-1}(r_{\Delta R_{ij}})R_j^TR_i\phi_i.
\end{align}
Comparing with a first-order Taylor expansion yields
\begin{equation}\label{eq:JresRi}
\boxed{\frac{\partial r_{\Delta R_{ij}}}{\partial\phi_i}
=-J_r^{-1}(r_{\Delta R_{ij}})R_j^TR_i.}
\end{equation}
The minus sign comes from
$(R_i\Exp(\phi_i))^T=\Exp(-\phi_i)R_i^T$;
the $R_j^TR_i$ term comes from the adjoint transformation.

\subsection{Derivative with respect to the final rotation}
Perturb $R_j$ while holding $R_i$ fixed:
\begin{align}
r_{\Delta R_{ij}}(R_j\Exp(\phi_j))
&=\Log\!\left(\Delta\tilde R_{ij}^{T}R_i^T
 R_j\Exp(\phi_j)\right)\notag\\
&=\Log\!\left(\Exp(r_{\Delta R_{ij}})
 \Exp(\phi_j)\right)\notag\\
&\simeq r_{\Delta R_{ij}}+J_r^{-1}(r_{\Delta R_{ij}})\phi_j.
\end{align}
Here the perturbation is already at the far right, so no adjoint
rearrangement is required. Therefore,
\begin{equation}\label{eq:JresRj}
\boxed{\frac{\partial r_{\Delta R_{ij}}}{\partial\phi_j}
=J_r^{-1}(r_{\Delta R_{ij}}).}
\end{equation}

\subsection{Summary of the two rotation-residual Jacobians}
\begin{align}
\frac{\partial r_{\Delta R_{ij}}}{\partial\phi_i}
&=-J_r^{-1}(r_{\Delta R_{ij}})R_j^TR_i,\\
\frac{\partial r_{\Delta R_{ij}}}{\partial\phi_j}
&=J_r^{-1}(r_{\Delta R_{ij}}).
\end{align}
Both are $3\times3$ Jacobians with respect to right perturbations.
These are the endpoint-rotation derivatives corresponding to the
textbook's (4.37)--(4.38). The following sections extend these results to the bias, velocity, and position residual Jacobians.


\section{Remaining IMU-Factor Residual Jacobians}
\subsection{Rotation residual with respect to gyroscope bias}
Let $b_{g,i}$ denote the original bias linearization point.
The optimizer has already applied the correction $\delta b_{g,i}$,
so its current bias estimate is
\begin{equation}
b_g'=b_{g,i}+\delta b_{g,i}.
\end{equation}
To compute the residual Jacobian, we perturb this \emph{current}
estimate by an additional infinitesimal increment $\epsilon_g$:
\begin{equation}
b_g''=b_g'+\epsilon_g
=b_{g,i}+\delta b_{g,i}+\epsilon_g.
\end{equation}
The important distinction is that the original pre-integration
and its bias Jacobian $\frac{\partial\Delta\tilde R_{ij}}{\partial b_{g,i}}$ were calculated at $b_{g,i}$.
We do \emph{not} recompute either of them at $b_g'$.
Instead, we differentiate the first-order bias-correction model
constructed at the original linearization point.
Retain $\frac{\partial\Delta\tilde R_{ij}}{\partial b_{g,i}}$ as defined earlier and set
\begin{equation}
J_{r,b}\defeq J_r\bigl(\frac{\partial\Delta\tilde R_{ij}}{\partial b_{g,i}}\delta b_{g,i}\bigr).
\end{equation}
The currently corrected measurement and residual are
\begin{align}
\Delta\tilde R'_{ij}&=\Delta\tilde R_{ij}
\Exp(\frac{\partial\Delta\tilde R_{ij}}{\partial b_{g,i}}\delta b_{g,i}),\\
r'_{\Delta R_{ij}}&=\Log\bigl((\Delta\tilde R'_{ij})^TR_i^TR_j\bigr).
\end{align}
Although $\epsilon_g$ perturbs the current estimate $b_g'$,
the fixed-linearization correction model still expands the total
bias offset $\delta b_{g,i}+\epsilon_g$ around $b_{g,i}$:
\begin{align}
\Delta\tilde R_{ij}(b_g' + \epsilon_g)
&\simeq\Delta\tilde R_{ij}(b_{g,i})
\Exp\bigl(\frac{\partial\Delta\tilde R_{ij}}{\partial b_{g,i}}(\delta b_{g,i}+\epsilon_g)\bigr).
\end{align}
Now use the first-order BCH formula with
$\phi=\frac{\partial\Delta\tilde R_{ij}}{\partial b_{g,i}}\delta b_{g,i}$ and
$\delta\phi=\frac{\partial\Delta\tilde R_{ij}}{\partial b_{g,i}}\epsilon_g$ to separate the new perturbation:
\begin{align}
&\Delta\tilde R_{ij}(b_{g,i}+\delta b_{g,i}+\epsilon_g)\notag\\
&\quad\simeq\Delta\tilde R_{ij}
\Exp\bigl(\frac{\partial\Delta\tilde R_{ij}}{\partial b_{g,i}}(\delta b_{g,i}+\epsilon_g)\bigr)\notag\\
&\quad\simeq\Delta\tilde R'_{ij}
\Exp(J_{r,b}\frac{\partial\Delta\tilde R_{ij}}{\partial b_{g,i}}\epsilon_g).
\end{align}
Here $J_{r,b}=J_r(\frac{\partial\Delta\tilde R_{ij}}{\partial b_{g,i}}\delta b_{g,i})$ is the right
Jacobian of the \emph{bias-induced rotation correction}.
The resulting local sensitivity $J_{r,b}\frac{\partial\Delta\tilde R_{ij}}{\partial b_{g,i}}$ is the
derivative of the \emph{fixed first-order correction model} at
the current correction $\delta b_{g,i}$. It need not equal the
Jacobian obtained by recomputing the original IMU pre-integration
at $b_g'$; the latter would account for higher-order bias dependence.

Substitute this into the rotation residual and use
$(AB)^T=B^TA^T$:
\begin{align}
r_{\Delta R_{ij}}(\epsilon_g)
&\simeq\Log\!\left(
\Exp(-J_{r,b}\frac{\partial\Delta\tilde R_{ij}}{\partial b_{g,i}}\epsilon_g)
(\Delta\tilde R'_{ij})^TR_i^TR_j\right)\notag\\
&=\Log\!\left(\Exp(-J_{r,b}\frac{\partial\Delta\tilde R_{ij}}{\partial b_{g,i}}\epsilon_g)
\Exp(r'_{\Delta R_{ij}})\right)\notag\\
&=\Log\!\left(\Exp(r'_{\Delta R_{ij}})
\Exp\bigl(-\Exp(r'_{\Delta R_{ij}})^T
J_{r,b}\frac{\partial\Delta\tilde R_{ij}}{\partial b_{g,i}}\epsilon_g\bigr)\right)\notag\\
&\simeq r'_{\Delta R_{ij}}
-J_r^{-1}(r'_{\Delta R_{ij}})
\Exp(r'_{\Delta R_{ij}})^T
J_{r,b}\frac{\partial\Delta\tilde R_{ij}}{\partial b_{g,i}}\epsilon_g.
\end{align}
The third line uses the exact $\SO(3)$ adjoint identity;
the last line uses the first-order right-BCH formula. Thus
\begin{equation}\label{eq:JresBg}
\boxed{\frac{\partial r_{\Delta R_{ij}}}{\partial\epsilon_g}
=-J_r^{-1}(r'_{\Delta R_{ij}})
\Exp(r'_{\Delta R_{ij}})^TJ_{r,b}\frac{\partial\Delta\tilde R_{ij}}{\partial b_{g,i}}.}
\end{equation}

\subsection{Velocity residual Jacobians}
Recall the velocity residual directly, without introducing an
additional shorthand for the state-predicted velocity:
\begin{equation}
r_{\Delta v_{ij}}=R_i^T(v_j-v_i-g\Delta t_{ij})-\Delta\tilde v_{ij}.
\end{equation}
Its Euclidean derivatives are
\begin{equation}
\boxed{\frac{\partial r_{\Delta v_{ij}}}{\partial v_i}=-R_i^T,
\qquad\frac{\partial r_{\Delta v_{ij}}}{\partial v_j}=R_i^T.}
\end{equation}
Using the right rotation perturbation $R_i'=R_i\Exp(\phi_i)$,
\begin{align}
r_{\Delta v_{ij}}(R_i')
&=\Exp(-\phi_i)R_i^T(v_j-v_i-g\Delta t_{ij})
-\Delta\tilde v_{ij}\notag\\
&\simeq r_{\Delta v_{ij}}+
\bigl[R_i^T(v_j-v_i-g\Delta t_{ij})\bigr]^\wedge\phi_i.
\end{align}
Consequently,
\begin{equation}\label{eq:JresVphi}
\boxed{\frac{\partial r_{\Delta v_{ij}}}{\partial\phi_i}
=\bigl[R_i^T(v_j-v_i-g\Delta t_{ij})\bigr]^\wedge.}
\end{equation}
Since the pre-integrated measurement is subtracted, its bias
Jacobians enter with negative signs:
\begin{equation}
\frac{\partial r_{\Delta v_{ij}}}{\partial b_{g,i}}
=-\frac{\partial\Delta\tilde v_{ij}}{\partial b_{g,i}},\qquad
\frac{\partial r_{\Delta v_{ij}}}{\partial b_{a,i}}
=-\frac{\partial\Delta\tilde v_{ij}}{\partial b_{a,i}}.
\end{equation}

\subsection{Position residual Jacobians}
The position residual is
\begin{equation}
r_{\Delta p_{ij}}=R_i^T\left(p_j-p_i-v_i\Delta t_{ij}-\tfrac12g\Delta t_{ij}^2\right)-\Delta\tilde p_{ij}.
\end{equation}
The Euclidean derivatives are
\begin{align}
\frac{\partial r_{\Delta p_{ij}}}{\partial p_i}&=-R_i^T,&
\frac{\partial r_{\Delta p_{ij}}}{\partial p_j}&=R_i^T,&
\frac{\partial r_{\Delta p_{ij}}}{\partial v_i}&=-R_i^T\Delta t_{ij}.
\end{align}
With $R_i'=R_i\Exp(\phi_i)$,
\begin{align}
r_{\Delta p_{ij}}(R_i')
&=\Exp(-\phi_i)R_i^T\left(p_j-p_i-v_i\Delta t_{ij}
-\tfrac12g\Delta t_{ij}^2\right)
-\Delta\tilde p_{ij}\notag\\
&\simeq r_{\Delta p_{ij}}+
\left[R_i^T\left(p_j-p_i-v_i\Delta t_{ij}
-\tfrac12g\Delta t_{ij}^2\right)\right]^\wedge\phi_i.
\end{align}
Therefore,
\begin{equation}\label{eq:JresPphi}
\boxed{\frac{\partial r_{\Delta p_{ij}}}{\partial\phi_i}
=\left[R_i^T\left(p_j-p_i-v_i\Delta t_{ij}
-\tfrac12g\Delta t_{ij}^2\right)\right]^\wedge.}
\end{equation}
The bias derivatives are
\begin{equation}
\frac{\partial r_{\Delta p_{ij}}}{\partial b_{g,i}}
=-\frac{\partial\Delta\tilde p_{ij}}{\partial b_{g,i}},\qquad
\frac{\partial r_{\Delta p_{ij}}}{\partial b_{a,i}}
=-\frac{\partial\Delta\tilde p_{ij}}{\partial b_{a,i}}.
\end{equation}

\section{Implementation Overview}
This section briefly connects the mathematical pre-integration model
with the graph-optimization implementation. It introduces no new
measurement model or Jacobian derivation.

\begin{enumerate}
  \item \textbf{Integrate IMU measurements.} For each IMU interval,
  accumulate $\Delta\tilde R_{ij}$, $\Delta\tilde v_{ij}$, and
  $\Delta\tilde p_{ij}$ using the fixed bias linearization point.
  \item \textbf{Propagate uncertainty and bias sensitivities.}
  Update the pre-integration covariance $\Sigma_{ij}$ and the five
  bias Jacobians alongside the accumulated measurements.
  \item \textbf{Construct the IMU factor.} Between keyframes $i$ and $j$,
  evaluate the nine-dimensional residual
  $r_{\mathrm{pre}}=(r_{\Delta R_{ij}}^T,
  r_{\Delta v_{ij}}^T,r_{\Delta p_{ij}}^T)^T$ and use
  $\Sigma_{ij}^{-1}$ as its information matrix, when invertible.
  Bias random-walk factors separately constrain the endpoint biases.
  \item \textbf{Optimize the navigation states.} Combine the IMU
  factor with any available additional observation factors and
  optimize the keyframe states using residuals and their Jacobians.
  Changes in the bias estimate are handled with the first-order
  bias correction model (or re-integration when necessary).
\end{enumerate}

In short, IMU pre-integration produces a \emph{relative-motion
measurement and its uncertainty}; the graph optimizer estimates the
\emph{navigation states}. This distinction parallels the separation
between propagation and measurement update in ESKF, although the
optimization mechanism is different.


\section{Conceptual Summary: From IMU Measurements to Graph Optimization}
\label{sec:conceptual_summary}
The main goal of IMU pre-integration is to compress many high-rate IMU
samples between two keyframes into one relative-motion constraint for
graph optimization. The pre-integration computation, measurement model,
and residual are not three separate estimation algorithms: they are
three connected stages of the same construction.

\subsection{Three distinct objects}
\begin{description}[leftmargin=2.6cm,style=nextline]
\item[Pre-integration computation]
Accumulates the IMU readings into
$\Delta\tilde R_{ij}$, $\Delta\tilde v_{ij}$, and $\Delta\tilde p_{ij}$.
These are measurements, not optimized navigation states.
\item[Measurement model]
Relates those accumulated measurements to the navigation states at
keyframes $i$ and $j$, including the effects of gravity and noise.
\item[Residual]
Measures the disagreement between the relative motion predicted by
the current navigation states and the pre-integrated measurements.
\end{description}

\subsection{Step 1: accumulate relative-motion measurements}
For each IMU sample $k$, the pre-integration recursively updates
$\Delta\tilde R_{i,k+1}$, $\Delta\tilde v_{i,k+1}$, and
$\Delta\tilde p_{i,k+1}$. For example,
\begin{equation}
\Delta\tilde R_{ij}=\prod_{k=i}^{j-1}
\Exp\bigl((\tilde\omega_k-b_{g,i})\Delta t_k\bigr).
\end{equation}
The corresponding velocity increment is
\begin{equation}
\Delta\tilde v_{ij}=\sum_{k=i}^{j-1}
\Delta\tilde R_{ik}(\tilde a_k-b_{a,i})\Delta t_k.
\end{equation}
The position increment is accumulated recursively using the preceding
velocity increment and the measured specific force. Crucially, these
calculations do not require the current world-frame estimates of
$R_i$, $p_i$, or $v_i$. Thus changing those state estimates during
optimization does not require re-integrating the raw IMU samples.

\subsection{Step 2: relate measurements to keyframe states}
The optimization variables at a keyframe are
\begin{equation}
x_k=(R_k,p_k,v_k,b_{a,k},b_{g,k}).
\end{equation}
The pre-integration measurement model provides the bridge between the
accumulated IMU quantities and these unknown states:
\begin{subequations}\label{eq:concept_measurement_model}
\begin{align}
\Delta\tilde R_{ij}
&\simeq R_i^TR_j\Exp(\delta\phi_{ij}),\\
\Delta\tilde v_{ij}
&=R_i^T(v_j-v_i-g\Delta t_{ij})+\delta v_{ij},\\
\Delta\tilde p_{ij}
&=R_i^T\!\left(p_j-p_i-v_i\Delta t_{ij}
-\tfrac12g\Delta t_{ij}^2\right)+\delta p_{ij}.
\end{align}
\end{subequations}
Here $\delta\phi_{ij}$, $\delta v_{ij}$, and $\delta p_{ij}$
are accumulated \emph{measurement-noise perturbations}, not optimization
increments applied to the keyframe states.

\subsection{Step 3: construct a residual}
The optimizer evaluates how well its current state estimates satisfy
the IMU measurement model. Using the residual convention in this chapter,
\begin{subequations}\label{eq:concept_residuals}
\begin{align}
r_{\Delta R_{ij}}
&=\Log\!\left(\Delta\tilde R_{ij}^TR_i^TR_j\right),\\
r_{\Delta v_{ij}}
&=R_i^T(v_j-v_i-g\Delta t_{ij})-\Delta\tilde v_{ij},\\
r_{\Delta p_{ij}}
&=R_i^T\!\left(p_j-p_i-v_i\Delta t_{ij}
-\tfrac12g\Delta t_{ij}^2\right)-\Delta\tilde p_{ij}.
\end{align}
\end{subequations}
Stacking these three three-dimensional residuals gives
\begin{equation}
r_{\mathrm{pre}}=
\begin{bmatrix}r_{\Delta R_{ij}}\\r_{\Delta v_{ij}}\\r_{\Delta p_{ij}}\end{bmatrix}
\in\R^9.
\end{equation}
For example, if the IMU-derived velocity increment is
$\Delta\tilde v_{ij}=(2,0,0)^T\,\mathrm{m/s}$ but the current states
predict $(1,0,0)^T\,\mathrm{m/s}$ (with $R_i=I$ and gravity removed),
then $r_{\Delta v_{ij}}=(-1,0,0)^T\,\mathrm{m/s}$. The optimizer
adjusts the states to reduce this discrepancy while satisfying other
measurement factors.

\subsection{Why covariance and bias Jacobians are needed}
\textbf{Pre-integration covariance.} IMU measurement noise accumulates
through the same integration process. Its first-order propagation is
\begin{equation}
\eta_{i,k+1}=A_k\eta_{ik}+B_k\eta_{d,k},\qquad
\Sigma_{i,k+1}=A_k\Sigma_{ik}A_k^T+B_kQ_{d,k}B_k^T.
\end{equation}
The resulting $\Sigma_{ij}\in\R^{9\times9}$ describes the
uncertainty of the relative-motion measurement, \emph{not} the
covariance of the optimized keyframe states. It supplies the
information matrix $\Omega_{ij}=\Sigma_{ij}^{-1}$ (when invertible)
for the weighted residual cost
\begin{equation}
E_{\mathrm{IMU}}=r_{\mathrm{pre}}^T\Omega_{ij}r_{\mathrm{pre}}.
\end{equation}

\textbf{Bias Jacobians.} Pre-integration uses fixed gyroscope and
accelerometer bias linearization values $b_{g,i}$ and $b_{a,i}$.
Because the graph optimizer also estimates the biases, changes in
these values alter the accumulated measurements. The five bias
Jacobians provide a first-order correction without re-integrating
every IMU sample at each iteration; sufficiently large changes may
still require re-integration.

\textbf{Residual Jacobians.} These are different from both the noise
propagation matrices and the bias-correction Jacobians: they describe
how $r_{\mathrm{pre}}$ changes when the \emph{optimization states} are
perturbed, and are used to construct the graph optimizer's linearized
least-squares system.

\subsection{The complete picture and comparison with ESKF}
\begin{center}
\fbox{\begin{minipage}{0.88\linewidth}
\centering
Raw IMU samples $\tilde\omega_k,\tilde a_k$\\[3pt]
$\Downarrow$\\[3pt]
Pre-integrated increments $\Delta\tilde R_{ij},\Delta\tilde v_{ij},\Delta\tilde p_{ij}$\\
Covariance $\Sigma_{ij}$ and five bias Jacobians\\[3pt]
$\Downarrow$\\[3pt]
Measurement model relating increments to $x_i,x_j$\\[3pt]
$\Downarrow$\\[3pt]
IMU factor residual $r_{\mathrm{pre}}\in\R^9$\\[3pt]
$\Downarrow$\\[3pt]
Graph optimization (IMU and other sensor factors)\\[3pt]
$\Downarrow$\\[3pt]
Optimized navigation states and biases
\end{minipage}}
\end{center}

In Chapter 3, ESKF uses IMU readings to propagate a navigation-state
estimate and then incorporates observations through a Kalman update.
Here, IMU pre-integration instead creates a relative-motion
\emph{measurement factor}; graph optimization adjusts the keyframe
states jointly to reduce the weighted residuals from IMU and other
sensors. The underlying IMU kinematics are closely related, but the
quantities propagated and the way measurements are fused differ.

\noindent\textbf{Four takeaways.}
Pre-integration computes relative-motion measurements; the measurement
model relates them to the unknown navigation states; the residual
quantifies their disagreement; and graph optimization adjusts the
states to minimize the weighted disagreements.

\appendix
\section{Complete IMU-Factor Jacobian Matrix}
Use the residual ordering $(r_{\Delta R},r_{\Delta v},r_{\Delta p})$
and optimization perturbation ordering
$(\phi_i,\delta p_i,\delta v_i,\epsilon_g,\epsilon_a,
\phi_j,\delta p_j,\delta v_j)$.
Each variable block has three components, so the assembled Jacobian
has dimensions $9\times24$. For the rotation blocks, retain the shorthand
\begin{equation}
W\defeq J_r^{-1}(r'),\qquad
G\defeq-W\Exp(r')^TJ_{r,b}\frac{\partial\Delta\tilde R_{ij}}{\partial b_{g,i}}.
\end{equation}
The rotation blocks below use the bias-corrected residual $r'$;
for the other two residuals the measurement bias Jacobians are
those derived in the previous section. We display the matrix in
three $3\times24$ block rows to avoid a page-width overflow.
\begin{equation}\label{eq:imu_factor_full_jacobian}
\boxed{\begin{gathered}
\displaystyle J_{ij}=\frac{\partial r_{ij}}{\partial\delta X}
=\begin{bmatrix}J_R\\J_v\\J_p\end{bmatrix},\\[2mm]
J_R=\begin{bmatrix}
-WR_j^TR_i&0&0&G&0&W&0&0
\end{bmatrix},\\[1mm]
J_v=\resizebox{.94\textwidth}{!}{$\begin{bmatrix}
\bigl[R_i^T(v_j-v_i-g\Delta t_{ij})\bigr]^\wedge&0&-R_i^T&-\frac{\partial\Delta\tilde v_{ij}}{\partial b_{g,i}}&-\frac{\partial\Delta\tilde v_{ij}}{\partial b_{a,i}}&0&0&R_i^T
\end{bmatrix}$},\\[1mm]
J_p=\resizebox{.94\textwidth}{!}{$\begin{bmatrix}
\bigl[R_i^T(p_j-p_i-v_i\Delta t_{ij}-\tfrac12g\Delta t_{ij}^2)\bigr]^\wedge&-R_i^T&-R_i^T\Delta t_{ij}&-\frac{\partial\Delta\tilde p_{ij}}{\partial b_{g,i}}&-\frac{\partial\Delta\tilde p_{ij}}{\partial b_{a,i}}&0&R_i^T&0
\end{bmatrix}$}.
\end{gathered}}
\end{equation}
Here every displayed block is $3\times3$. The position and
velocity rows do not directly depend on $R_j$ in this residual
convention. The Jacobian above is evaluated at the current bias-corrected pre-integration. The velocity and position bias derivatives denote the stored first-order sensitivities.

\subsection{Scope and next steps}
The present notes stop at the position residual
Jacobians. The assembled matrix is a convenient grouping
of the individual Jacobian blocks, rather than a separate equation
in the textbook. No additional sensor factors are included here.


\subsection{Full two-keyframe state ordering}
The displayed $9\times24$ matrix uses only the eight three-dimensional
state blocks directly connected to the pre-integration factor.
For the complete two-keyframe ordering
\begin{equation}
 \delta X_{ij}^{\mathrm{full}}=
 (\phi_i,\delta p_i,\delta v_i,\delta b_{a,i},\delta b_{g,i},
 \phi_j,\delta p_j,\delta v_j,\delta b_{a,j},\delta b_{g,j}),
\end{equation}
the full Jacobian has size $9\times30$. It is obtained by
reordering the existing blocks to match this ordering and appending
two zero $9\times3$ blocks for the endpoint biases. Those endpoint
biases are constrained by the separate bias random-walk factors.

\end{document}
